\section{The signed divisor as an integer positive-definite function}
\label{sec:s5-crystal}

The sign condition at the good primes survives the finitely many bad
Euler factors in a useful precise form. The correction for those
factors is continuous on the zero side, and therefore disappears
when individual zero atoms are extracted. The topology in the
conclusion is essential: the resulting function is positive definite
on the additive group $\mathbb R_d$ with its discrete topology.

\paragraph{MA-FQ-001 (finite Euler corrections preserve atomic positivity).}
Retain the family, the specified comparison representation $\varpi$,
and $(\mathcal R_B)$ of MA-S5-001. Let
$\mathcal F_R=(f_R,K_R,m_R)$ be the difference of the actual
automorphic and Artin data in MA-LD-001, with its actual completed
divisor. Then
\[
 K_R=0,\qquad
 \mu=\sum_{t\in\mathbb R}m_R(t)\delta_t
       \ \text{is a tempered, locally finite signed measure},
 \qquad m_R(0)=-k.
 \tag{MA-FQ.1a}
\]
All its masses are integers. Define $a(t)=-m_R(t)$, with value zero
away from the divisor. The function $a$ is positive definite on
$\mathbb R_d$, and
\[
 a(0)=k,\qquad |m_R(t)|\le k\quad(t\in\mathbb R),\qquad
 a\in B(\mathbb R_d),\qquad \|a\|_{B(\mathbb R_d)}=k.
 \tag{MA-FQ.1b}
\]
Here $B(\mathbb R_d)$ is the Fourier--Stieltjes algebra: the
transforms of finite complex measures on the compact character
group $\widehat{\mathbb R_d}$. In particular $a$ has finite integer
range and locally finite support in the usual real topology.
No positivity of the signed divisor as a measure is asserted.

\begin{proof}
\emph{The exact archimedean parameters.}
The quadratic field $E=L^{A_5}$ is totally real. The finite-image
two-dimensional representation $\varrho$ used by Calegari is odd
at each real place of $E$, so its local parameter is
$\mathbf1\oplus\operatorname{sgn}$. Calegari's construction gives
\[
 \operatorname{BC}_{E/\mathbb Q}(\varpi)
  =\operatorname{Sym}^4(\pi_E)\otimes\det(\pi_E)^{-2},
 \tag{MA-FQ.1c}
\]
where $\pi_E$ is the automorphic representation corresponding to
$\varrho$. At a real place the right side has parameter
$\mathbf1^{\oplus3}\oplus\operatorname{sgn}^{\oplus2}$:
the five symmetric-power monomials have signs $+,-,+,-,+$,
and the determinant twist has sign $\operatorname{sgn}^{-2}=1$.
Every real place of $\mathbb Q$ splits in the real quadratic
extension. Local base change there is the identity on the real
parameter, so the same parameter belongs to $\varpi_\infty$.
There is no descent ambiguity at infinity; in particular
$\eta_\infty=1$.

For $W$, complex conjugation has order two and trace one in
dimension five. Its $+1$ and $-1$ multiplicities are therefore
three and two. Thus the two actual gamma factors are exactly
\[
 G_{\varpi}(s)=G_W(s)
  =\Gamma_{\mathbb R}(s)^3\Gamma_{\mathbb R}(s+1)^2
  =\Gamma_{\mathbb R}(s)\Gamma_{\mathbb C}(s)^2.
 \tag{MA-FQ.1d}
\]
Booker's formula for the kernel in Example~1.4(2), also used for
Artin representations in Example~1.4(3), consequently gives
$K_R=0$. This proves equality of the kernels, not merely equality
of their residues at zero. The sources are Calegari's character
table, author-hosted publisher PDF p.~4; Lemma~3.2 and its
preceding paragraph, p.~8; and the opening of \S4, pp.~8--9.
The invoked modularity, symmetric-power transfer and cyclic
base-change results remain the classical inputs used in
Calegari's construction. This deduction does not assume the
desired all-place comparison over $\mathbb Q$.

\emph{The divisor and the precise Fourier identity.}
MA-LD-001 proves that the complete signed divisor of this actual
quotient is supported on $\operatorname{Re}s=1/2$, using
$(\mathcal R_B)$ and MA-S5-004; if $k=0$ the assertion follows
already from MA-S5-001. Multiplication by the nonzero conductor
exponential changes no order. Hence $m_R$ is integer valued,
real supported, and has $m_R(0)=-k$. Axiom~(A3) of Booker's
Definition~1.3 gives constants $C,A$ such that
\[
 \sum_{|t|\le T}|m_R(t)|\le C(1+T)^A.
 \tag{MA-FQ.1e}
\]
Integrality implies that there are only finitely many nonzero
masses in each compact interval. Summation over dyadic annuli
shows that $\int(1+|t|)^{-N}\,d|\mu|(t)<\infty$ for $N>A$.
Thus $\mu$ defines a tempered distribution as well as a locally
finite signed measure.

Use the following Fourier convention throughout this argument:
\[
 \widehat u(x)=\int_{\mathbb R}e^{itx}u(t)\,dt,
 \qquad
 u(t)=\frac1{2\pi}\int_{\mathbb R}e^{-itx}\widehat u(x)\,dx,
 \qquad
 \langle\widehat T,g\rangle=\langle T,\widehat g\rangle.
 \tag{MA-FQ.1f}
\]
Let $\mathcal B$ be the finite bad-prime set, disjoint from
$\Omega$. Define the positive locally finite measure
\[
 M_\Omega=\sum_{\substack{p\in\Omega\\j\ge1\ \mathrm{odd}}}
 \frac{4\log p}{p^{j/2}}
   \bigl(\delta_{j\log p}+\delta_{-j\log p}\bigr)
 \tag{MA-FQ.1g}
\]
and the possibly signed or complex measure
\[
 B_0=2f_R(1)\delta_0+
 \sum_{p\in\mathcal B}\sum_{j\ge1}
 \left(f_R(p^j)\delta_{j\log p}
       +\overline{f_R(p^j)}\delta_{-j\log p}\right).
 \tag{MA-FQ.1h}
\]
Its total variation is finite. Indeed, axiom~(A1) with exponent
two gives
\[
 \sum_{j\ge1}|f_R(p^j)|
 \le C_2\sum_{j\ge1}(j\log p)^{-2}<\infty
 \quad\text{for each fixed }p,
 \tag{MA-FQ.1i}
\]
and $\mathcal B$ is finite. This argument uses no sign assumption
on any bad Euler coefficient. The stronger root-of-unity
description of Calegari's bad local polynomials in Lemma~4.8
is consistent with this bound but is not needed here.

The measure $M_\Omega$ is tempered. For $k>0$, MA-S5-004 gives
$\sum_{p\le e^T,\ p\in\Omega}(\log p)/\sqrt p=O(1+T)$.
The total mass contributed by $j\ge3$ is bounded by
\[
 8\sum_p(\log p)\sum_{j\ge3}p^{-j/2}<\infty.
\]
Consequently $M_\Omega([-T,T])=O(1+T)$. If $k=0$, then
$\Omega$ is finite by MA-S5-001, and the full geometric tails
have finite mass, which gives the same conclusion.

For a compact smooth test $g$ satisfying
$g(-x)=\overline{g(x)}$, its Fourier transform in (MA-FQ.1f)
is real on $\mathbb R$. Booker's (A4), with $K_R=0$, reads
\[
 \langle\mu,\widehat g\rangle
 =-2\operatorname{Re}\sum_{n\ge1}f_R(n)g(\log n).
 \tag{MA-FQ.1j}
\]
The summation includes $n=1$ under the minus sign. Since
$f_R(1)$ and $g(0)$ are real, the term at zero is
$-2f_R(1)g(0)$. For $n\ge2$, the two atoms in (MA-FQ.1h)
give exactly twice the real part of $f_R(n)g(\log n)$.
The coefficients at good primes are the real nonnegative
coefficients in (MA-LD.6). All other coefficients with $n>1$
vanish unless $n$ is a prime power. It follows that
\[
                  \widehat{(-\mu)}=M_\Omega+B_0.
 \tag{MA-FQ.1k}
\]
To see that this is a complex distribution identity, decompose
an arbitrary compact smooth $g$ as $g=g_1+ig_2$, where
\[
 g_1(x)=\frac{g(x)+\overline{g(-x)}}2,
 \qquad
 g_2(x)=\frac{g(x)-\overline{g(-x)}}{2i}.
\]
Both $g_1$ and $g_2$ have the required Hermitian symmetry, and
both sides of (MA-FQ.1k) are complex linear in the test.
The identity on compact smooth tests extends to Schwartz tests
because all three distributions are tempered and compact
smooth functions are dense in the Schwartz space.

\emph{Cancellation by a continuous density.}
Define
\[
 q(t)=-\frac1{2\pi}\int_{\mathbb R}e^{-itx}\,dB_0(x),
 \qquad \nu=-\mu+q(t)\,dt.
 \tag{MA-FQ.1l}
\]
Finite total variation makes $q$ bounded and continuous, with
$\|q\|_\infty\le\|B_0\|_{\mathrm{TV}}/(2\pi)$.
Fourier inversion and Fubini on Schwartz tests give
$\widehat{q(t)\,dt}=-B_0$, so
$\widehat\nu=M_\Omega$. In particular $\nu$ is a distribution
of positive type. More explicitly, with
$\widetilde\psi(t)=\overline{\psi(-t)}$,
\[
 \langle\nu,\psi*\widetilde\psi\rangle
   =\frac1{2\pi}\int_{\mathbb R}
                         |\widehat\psi(-x)|^2\,dM_\Omega(x)
   \ge0\qquad(\psi\in C_c^\infty(\mathbb R)).
 \tag{MA-FQ.1m}
\]
The reflection $-x$ and the factor $(2\pi)^{-1}$ follow from
the Fourier convention (MA-FQ.1f). They do not affect positivity.

\emph{Extraction of the atomic coefficient function.}
Choose a nonzero real even $\phi\in C_c^\infty(\mathbb R)$.
Put $\phi_\epsilon(t)=\epsilon^{-1/2}\phi(t/\epsilon)$ and
$\Psi=\phi*\widetilde\phi$. Then
$\phi_\epsilon*\widetilde\phi_\epsilon(t)=\Psi(t/\epsilon)$,
where $\Psi$ is compactly supported and
$\Psi(0)=\|\phi\|_2^2>0$. For arbitrary finite choices
$x_1,\ldots,x_r\in\mathbb R$ and $c_1,\ldots,c_r\in\mathbb C$,
apply (MA-FQ.1m) to
$\psi_\epsilon=\sum_i c_i\phi_\epsilon(\,\cdot-x_i)$.
This gives
\[
 0\le\sum_{i,j}c_i\overline{c_j}\,
 \left\langle\nu,
  \Psi\left(\frac{\,\cdot-(x_i-x_j)}\epsilon\right)
 \right\rangle.
 \tag{MA-FQ.1n}
\]
There are only finitely many differences $x_i-x_j$.
Local finiteness of the support of $\mu$ therefore allows a
single sufficiently small $\epsilon$ for which the atomic
contribution at each such difference is precisely
$a(x_i-x_j)\Psi(0)$; there are no other atoms in the
corresponding shrinking support. The continuous contribution
has absolute value at most
$\epsilon\|q\|_\infty\|\Psi\|_1$ for each pair $i,j$.
Letting $\epsilon$ decrease to zero in (MA-FQ.1n) and dividing
by $\Psi(0)$ gives
\[
       \sum_{i,j}c_i\overline{c_j}\,a(x_i-x_j)\ge0.
 \tag{MA-FQ.1o}
\]
This is exactly positive definiteness on $\mathbb R_d$.
The two-by-two principal matrices give
$a(-t)=\overline{a(t)}$ and $|a(t)|\le a(0)=k$.

For completeness the standard Bochner implication here requires
no continuity in the usual real topology. Form the pre-Hilbert
space with formal vectors $e_x$, inner products
$\langle e_x,e_y\rangle=a(x-y)$, and quotient its nullspace.
Its completion carries commuting unitary translations
$U_z e_x=e_{x+z}$ and a vector $v=e_0$ of squared norm $k$.
If $k=0$ all the coefficients vanish. Otherwise the unital
commutative $C^*$-algebra generated by the $U_z$ is $C(X)$
for its compact spectrum $X$. A point $\omega\in X$ determines
the character $z\mapsto\omega(U_z)$ of $\mathbb R_d$.
The positive functional $T\mapsto\langle Tv,v\rangle$ is,
by the Riesz representation theorem, integration against a
positive measure of mass $k$ on $X$. Push it forward to
$\widehat{\mathbb R_d}$ to obtain a positive finite measure
$\sigma$ with
\[
        a(z)=\int_{\widehat{\mathbb R_d}}\chi(z)\,d\sigma(\chi),
        \qquad \sigma(\widehat{\mathbb R_d})=k.
 \tag{MA-FQ.1p}
\]
The map to the character group is continuous in its product
topology, so the pushforward is well defined even though the
compact character group need not be metrizable. This proves
$a\in B(\mathbb R_d)$ with norm at most $k$. Every representing
measure has total variation at least $|a(0)|=k$, proving the
norm equality. Finally $a$ is integer valued, so the bound
$|a|\le k$ gives finite range. Its support was already shown
to be locally finite. This proves every assertion.
\end{proof}

\paragraph{Proof scope and cited input.}
The argument above proves the atomic positivity assertion even
with arbitrary signs in the finitely many bad-prime progressions.
It does not require that $-\mu$ itself be of positive type;
the continuous correction in (MA-FQ.1l) is part of the proof.
No theorem about Fourier quasicrystals with uniformly discrete
support is invoked: uniform discreteness has not been assumed.
The next coset-ring step uses the locally compact group
$\mathbb R_d$, while local finiteness always uses the usual
real topology.

The exact Booker locators are Definition~1.3(A1), (A3), (A4),
author p.~3, and Examples~1.4(2)--(3), p.~4, in
\path{arXiv:1308.3067v2}. The actual critical-line-support
argument and its complete source-proof audit are MA-LD-001.
The compact-group Fourier--Stieltjes convention is the one used
in Sanders (2020), \S12, before Theorem~12.1; replacing each
character by its inverse changes neither the algebra nor its norm.
The last paragraph of the proof supplies the required special
case of the usual Bochner argument from the commutative
$C^*$-algebra and Riesz representation theorems, which are
classical background inputs.

Calegari, F. (2013). The Artin conjecture for some $S_5$-extensions.
\emph{Mathematische Annalen, 356}, 191--207.
\url{https://doi.org/10.1007/s00208-012-0839-4}.

Booker, A. R. (2015). $L$-functions as distributions.
\emph{Mathematische Annalen, 363}, 423--454.
\url{https://doi.org/10.1007/s00208-015-1178-z}.
